\chapter{Acotación total y compacidad local}\label{chap:15}

\section{Qué le falta a un conjunto cerrado y acotado}
En $\R^n$, ser cerrado y acotado equivale a ser compacto. En $\ell^\infty$ ya encontramos una bola cerrada que contiene infinitos puntos separados por una distancia fija. La noción que excluye esta situación es la acotación total: a cualquier escala, el espacio debe poder describirse con finitos centros. La completitud garantiza, además, que los límites que esa descripción produce permanezcan dentro del espacio. Usaremos compacidad secuencial y completitud. Desarrollamos la lectura 15 de \cite{valencia}.

\section{La caracterización métrica de la compacidad}
\begin{definicion}
Un espacio métrico $X$ es \emph{totalmente acotado}\index{acotacion total@acotación total} si para cada $\eps>0$ existen $x_1,\ldots,x_m\in X$ tales que $X=\bigcup_{i=1}^mB(x_i,\eps)$. Los centros constituyen una \emph{red finita de radio $\eps$}. Aquí ``red'' significa conjunto de centros, no una familia indexada por un conjunto dirigido.
\end{definicion}
Un espacio totalmente acotado no vacío es acotado: basta una cubierta por finitas bolas de radio $1$ y estimar las distancias de los centros a uno de ellos. Todo subconjunto $A$ de un espacio totalmente acotado es totalmente acotado con su métrica relativa. Para probarlo, se cubre el espacio con bolas de radio $\eps/2$ y, de cada bola que toca $A$, se elige un centro nuevo en $A$; las bolas de radio $\eps$ con los nuevos centros cubren $A$. Este cambio de centros evita suponer que los centros originales pertenecen a $A$.

\begin{teorema}\label{ch15:criterio}
Un espacio métrico es compacto si y solo si es completo y totalmente acotado.
\end{teorema}
\begin{proof}
Si es compacto, es completo por el ejercicio 3 del capítulo \ref{chap:13}. La cubierta por todas las bolas de un radio fijo tiene una subcubierta finita, lo que prueba la acotación total.

Para el recíproco, sea $(x_n)$ una sucesión. Una cubierta finita por bolas de radio $1/2$ contiene una bola a la que pertenecen infinitos términos: de lo contrario, la unión de las cantidades finitas de índices sería finita. Sea $I_1$ un conjunto infinito de esos índices. Cubrimos ahora $X$ por finitas bolas de radio $1/4$ y elegimos un conjunto infinito $I_2\subset I_1$ cuyos términos estén en una sola bola. Continuamos con radios $2^{-j}$. Elegimos índices $n_j\in I_j$ con $n_{j+1}>n_j$; esto es posible porque cada $I_j$ es infinito. Si $k,l\ge j$, entonces $n_k,n_l\in I_j$ y
\[
d(x_{n_k},x_{n_l})<2\cdot2^{-j}.
\]
La subsucesión es de Cauchy. La completitud le proporciona un límite en $X$. Así toda sucesión admite una subsucesión convergente y el capítulo \ref{chap:14} implica compacidad.
\end{proof}

\begin{observacion}[Heine--Borel y su alcance]
En $\R^n$, un conjunto acotado está contenido en una caja $[-M,M]^n$. Para cada $\eps>0$, dividimos cada intervalo en finitos subintervalos de longitud menor que $\eps/\sqrt n$. Cada caja pequeña tiene diámetro menor que $\eps$; eligiendo un punto del conjunto en cada caja que lo toque obtenemos una cubierta finita por bolas de radio $\eps$. Por tanto el conjunto es totalmente acotado. Si además es cerrado, es completo como subespacio cerrado del espacio completo $\R^n$, y el teorema da compacidad. La necesidad de cerrado y acotado ya se probó. En dimensión infinita, acotado no implica totalmente acotado, como muestran los vectores $e_n$ de la bola de $\ell^\infty$.\index{teorema de Heine--Borel}
\end{observacion}

\begin{proposicion}\label{ch15:relativa}
Para $A\subset X$ son equivalentes:
\begin{enumerate}[label=(\roman*)]
\item la cerradura $\cl A^{\,X}$ es compacta;
\item toda sucesión en $A$ tiene una subsucesión convergente a algún punto de $X$;
\item $A$ es totalmente acotado y $\cl A^{\,X}$ es completo.
\end{enumerate}
Cuando se cumplen, $A$ es \emph{relativamente compacto en $X$}.\index{compacidad!relativa}
\end{proposicion}
\begin{proof}
De (i) se sigue (ii) aplicando compacidad secuencial a la cerradura. También se sigue (iii), porque la cerradura es completa y totalmente acotada, y sus subconjuntos son totalmente acotados.

Para (ii)$\Rightarrow$(i), sea $(z_n)$ una sucesión en $\cl A$. Si $A$ es vacío no hay nada que probar. En otro caso, la definición de cerradura permite elegir $a_n\in A$ con $d(z_n,a_n)<1/n$. Por (ii), una subsucesión $a_{n_j}$ converge a $a\in X$. Entonces
\[
d(z_{n_j},a)\le 1/n_j+d(a_{n_j},a)\longrightarrow0.
\]
Como $\cl A$ es cerrado, $a\in\cl A$. Así $\cl A$ es secuencialmente compacto. Para (iii)$\Rightarrow$(i), el ejercicio 1 demuestra que $\cl A$ es totalmente acotado; se aplica el teorema \ref{ch15:criterio}.
\end{proof}

\begin{ejemplo}[El espacio ambiente importa]
El intervalo $(0,1)$ es relativamente compacto en $\R$, pues su cerradura allí es $[0,1]$. No es relativamente compacto en sí mismo: su cerradura relativa es $(0,1)$ y la sucesión $1/n$ no tiene ninguna subsucesión que converja a un punto de ese espacio. En un ambiente completo, ser relativamente compacto equivale a ser totalmente acotado. En un ambiente incompleto hay que comprobar también la completitud de la cerradura.
\end{ejemplo}

\section{Compacidad alrededor de cada punto}
\begin{definicion}
Un espacio métrico $X$ es \emph{localmente compacto}\index{compacidad!local} si cada punto tiene una vecindad compacta: para cada $x$ existen un abierto $O$ y un compacto $K$ con $x\in O\subset K$. Equivalentemente, cada punto tiene una vecindad cuya cerradura es compacta. Una vecindad no tiene por qué ser abierta.
\end{definicion}
La equivalencia se verifica así: una vecindad compacta es cerrada, por lo que es su propia cerradura; si una vecindad tiene cerradura compacta, esa cerradura sigue siendo una vecindad del punto.

\begin{proposicion}\label{ch15:subespacios}
Todo subespacio abierto o cerrado de un espacio métrico localmente compacto es localmente compacto.
\end{proposicion}
\begin{proof}
Sea $x\in A$ y sea $K$ una vecindad compacta de $x$ en $X$. Si $A$ es cerrado, $K\cap A$ es compacto y contiene una vecindad relativa de $x$ en $A$.

Si $A$ es abierto, elegimos un abierto $O$ con $x\in O\subset K$. Existe $r>0$ con $\overline B(x,r)\subset A\cap O$: basta escoger primero una bola de radio $2r$ contenida en ese abierto. La bola cerrada es un cerrado de $K$, luego es compacta, y constituye una vecindad de $x$ tanto en $X$ como en $A$.
\end{proof}

\begin{proposicion}\label{ch15:engrosamiento}
Si $X$ es localmente compacto y $A\subset X$ es compacto, existe $r>0$ tal que $\cl{B(A,r)}$ es compacto, donde $B(A,r)=\{x:\dist(x,A)<r\}$.
\end{proposicion}
\begin{proof}
Para $A=\varnothing$, interpretamos $B(A,r)=\varnothing$. En otro caso, por el ejercicio 3 cada $a\in A$ admite $r_a>0$ con $\overline B(a,r_a)$ compacta. Elegimos una subcubierta finita $B(a_i,r_i/2)$ de $A$, ponemos $r=\frac14\min_i r_i$ y $K=\bigcup_i\overline B(a_i,r_i)$. El conjunto $K$ es compacto. Si $z\in B(A,r)$, hay $a\in A$ con $d(z,a)<r$, y después un $i$ con $d(a,a_i)<r_i/2$. Así $d(z,a_i)<r+r_i/2<r_i$ y $z\in K$. Como $K$ es cerrado, también $\cl{B(A,r)}\subset K$. Esta cerradura es un cerrado de un compacto, luego compacta. No se ha supuesto que todas las bolas cerradas de $X$ sean compactas: solo se emplean radios locales adecuados.
\end{proof}

\section{Ejercicios resueltos}
\begin{ejercicio}[Lectura 15, ejercicio 1]
Si $A\subset X$ es totalmente acotado, probar que su cerradura también lo es.
\end{ejercicio}
\begin{solucion}
Fijado $\eps>0$, cubrimos $A$ con finitas bolas $B(a_i,\eps/2)$. Por ser la unión finita, su cerradura está contenida en la unión de las cerraduras de esas bolas, y cada una de estas está contenida en $\overline B(a_i,\eps/2)\subset B(a_i,\eps)$. En detalle, si un punto estuviera fuera de todas las cerraduras, la intersección de las finitas vecindades que evitan cada bola sería una vecindad que evita su unión. Por tanto las bolas $B(a_i,\eps)$ cubren $\cl A$.
\end{solucion}

\begin{ejercicio}[Lectura 15, ejercicio 2]
Probar que el recorrido de toda sucesión de Cauchy es totalmente acotado.
\end{ejercicio}
\begin{solucion}
Para $\eps>0$, elegimos $N$ con $d(x_n,x_m)<\eps$ si $n,m\ge N$. La cola está en $B(x_N,\eps)$, y los términos anteriores quedan cubiertos por $B(x_1,\eps),\ldots,B(x_{N-1},\eps)$. Esta es una cubierta finita del recorrido. No necesitamos que la sucesión converja ni que el ambiente sea completo.
\end{solucion}

\begin{ejercicio}[Lectura 15, ejercicio 3]
Probar que $X$ es localmente compacto si y solo si cada $x$ admite $r_x>0$ con la bola cerrada $\overline B(x,r_x)$ compacta.
\end{ejercicio}
\begin{solucion}
Si $K$ es una vecindad compacta de $x$, existe $s>0$ con $B(x,s)\subset K$. Tomamos $r_x=s/2$; la bola cerrada de ese radio es un cerrado de $K$, luego compacta. Para la inversa, $\overline B(x,r_x)$ contiene la bola abierta $B(x,r_x)$ y por tanto es una vecindad compacta. Se trata de la bola cerrada, no de afirmar que la bola abierta sea compacta.
\end{solucion}

\begin{ejercicio}[Lectura 15, ejercicio 4]
Probar que $\Q$, con la métrica inducida de $\R$, no es localmente compacto.
\end{ejercicio}
\begin{solucion}
Supongamos que $q\in\Q$ tiene una vecindad compacta $K\subset\Q$. Entonces $\Q\cap(q-r,q+r)\subset K$ para algún $r>0$. Elegimos un irracional $\alpha\in(q-r,q+r)$ y racionales $q_n$ en ese intervalo que converjan en $\R$ a $\alpha$. Todos pertenecen a $K$. Una subsucesión debería converger en $\Q$ a algún $a\in K$; la inclusión continua $\Q\hookrightarrow\R$ daría convergencia a $a$ en $\R$, contradiciendo la unicidad del límite $\alpha\notin\Q$. Así una subparte arbitraria de un espacio localmente compacto puede dejar de ser localmente compacta. Las hipótesis abierto o cerrado de la proposición \ref{ch15:subespacios} son suficientes; no se afirma que sean individualmente necesarias para cada subespacio.
\end{solucion}

\begin{ejercicio}[Lectura 15, ejercicio 5]
Sean $X$ localmente compacto, $A\subset X$ compacto y $U\subset X$ abierto con $A\subset U$. Probar:
\begin{enumerate}[label=(\roman*)]
\item existen $V$ abierto y $B$ compacto con $A\subset V\subset B\subset U$;
\item existe $f:X\to[0,1]$ continua, igual a $1$ en $A$, cuyo soporte es compacto y está contenido en $U$.
\end{enumerate}
\end{ejercicio}
\begin{solucion}
Para (i), cada $a\in A$ admite una bola cerrada compacta de radio $s_a>0$. Como $U$ es abierto, podemos disminuir el radio hasta elegir $r_a>0$ con $\overline B(a,2r_a)\subset U$ y con esa bola todavía compacta, por ser cerrada en la bola original. Extraemos una subcubierta finita $B(a_i,r_i)$ de $A$ y ponemos
\[
V=\bigcup_iB(a_i,r_i),\qquad B=\bigcup_i\overline B(a_i,r_i).
\]
Entonces $\cl V\subset B\subset U$ y $B$ es compacto.

Para (ii), si $A=\varnothing$, elegimos $f=0$. Para $A\ne\varnothing$ y $X\setminus V\ne\varnothing$, definimos
\[
f(x)=\frac{\dist(x,X\setminus V)}{\dist(x,X\setminus V)+\dist(x,A)}.
\]
Ambas distancias son continuas. El denominador nunca es cero: si $\dist(x,A)=0$, el cerrado $A$ contiene $x$, y entonces $x\in V$ tiene distancia positiva al cerrado $X\setminus V$. La función toma valores en $[0,1]$, vale $1$ en $A$ y vale $0$ fuera de $V$. Por tanto $\supp f=\cl{\{x:f(x)\ne0\}}\subset\cl V\subset B\subset U$; es cerrado en el compacto $B$, luego compacto. Si $V=X$, las inclusiones de (i) implican $B=X$ compacto y podemos elegir $f=1$. Esta construcción explica por qué se debe reservar espacio entre $A$ y el exterior de $U$.\index{funcion de corte@función de corte}\index{soporte}
\end{solucion}

\begin{ejercicio}[Lectura 15, ejercicio 6]
Si el compacto $A$ está cubierto por abiertos $U_1,\ldots,U_n$ de un espacio localmente compacto $X$, construir funciones continuas $f_k:X\to[0,1]$ de soporte compacto contenido en $U_k$, tales que $\sum_k f_k\le1$ en $X$ y $\sum_kf_k=1$ en $A$.
\end{ejercicio}
\begin{solucion}
Si $A$ es vacío, tomamos todas las funciones nulas. Para cada $a\in A$, elegimos $k(a)$ con $a\in U_{k(a)}$ y un abierto $V_a$ que contenga $a$ y tenga cerradura compacta contenida en $U_{k(a)}$, mediante el ejercicio 5(i) aplicado a $\{a\}$. El inciso (ii), aplicado al compacto $\cl V_a$, produce $h_a:X\to[0,1]$ igual a $1$ en $\cl V_a$, con soporte compacto contenido en $U_{k(a)}$. Elegimos $a_1,\ldots,a_m$ de modo que los $V_{a_j}$ cubran $A$.

Agrupamos por el abierto al que están subordinadas:
\[
g_k=\sum_{j:k(a_j)=k}h_{a_j},\qquad S=\sum_{k=1}^n g_k,\qquad
f_k=\frac{g_k}{\max\{1,S\}}.
\]
Una suma sin términos se interpreta como cero. El denominador es continuo y al menos $1$. Como $0\le g_k\le S\le\max\{1,S\}$, tenemos $0\le f_k\le1$. Su soporte está contenido en la unión finita de los soportes de los $h_{a_j}$ correspondientes; esta unión es compacta y está contenida en $U_k$. Finalmente,
\[
\sum_k f_k=\frac S{\max\{1,S\}}\le1.
\]
Si $x\in A$, algún $V_{a_j}$ contiene $x$, así que $h_{a_j}(x)=1$ y $S(x)\ge1$. Allí la suma vale exactamente $1$. Normalizar por $S$ sin el máximo no serviría en los puntos donde $S=0$.\index{particion de la unidad@partición de la unidad}
\end{solucion}

\section{Idea que conviene conservar}
La acotación total controla todas las escalas; la completitud retiene sus límites. La compacidad local permite construir funciones de corte y particiones finitas de la unidad, herramientas que más adelante harán posible localizar problemas de integración y aproximación.
